\chapter{Non-vacuity}\label{ch:nonvacuity}

A definition is non-vacuous when it is inhabited by the object leanVM intends and refuted by
the near misses; a theorem, when its hypotheses can be met by the real protocol and not only by
a trivial one, and when what it concludes is not already true of everything. This chapter
records what the review established by probe, in Lean, against \code{main} at \code{b435631},
and what it established on paper. The probes are reproduced in \cref{ch:probes}; the rule of the
review was that a conclusion resting on a probe names it.

\section{The relation and the seams}\label{sec:nv-relation}

\paragraph{\lean{M3Holds}.} On the toy instance the honest stack satisfies it, and for each of
the four checkable clauses a stack fails that clause alone: the repository's own tests do this
(\file{tests/LeanerVMTests/Protocol/Spine.lean}). The review added the cases those tests lack
(probe \code{P2Relation}): a tuple pushed twice and pulled once fails \lean{Balanced}, and a
tuple pushed twice and never pulled fails it too, which a balance summed in the field would
accept; the side filter and the row ranges of \lean{tuples} were exercised on instances built
for the purpose. Clause by clause against the pinned specification, \lean{M3Holds} is the
relation the bus, the table sumcheck, the public-input check and Flock together establish
(\cref{ch:faithfulness}); the review found no clause wrong.

\paragraph{The seams.} Each seam was inhabited and refuted per conjunct (probe \code{P3aSeams});
the zerocheck escape, a constraint violated on the cube whose extension vanishes at the point
$(1)$, was exhibited on a stack \code{badConstraint} to show that the bus seam's claims can hold
of a stack outside \lean{M3Holds}, which is what the error of the bus phase pays for.
\lean{Seam.done} is \lean{True}, so the last law of a knowledge state function reads
\enquote{if the verifier accepts then the state is true} (probe \code{NonVacuity}); this is
what an output relation carrying nothing should give.

\paragraph{What the seams admit beyond leanVM.} \lean{Seam.bus} admits statements the deployed
bus phase never emits and the deployed table sumcheck cannot serve: a statement true of an
honest stack yet outside the seam by its degree clause alone; a linear claim with two terms of
one table at two points; a thousand copies of a true claim (probes \code{SeamBusShape},
\code{SeamBusMember}). This is not a vacuity of the seam but a burden on the phase that
consumes it, developed in \cref{sec:faith-table}.

\section{The hypotheses of the master theorems}\label{sec:nv-hypotheses}

The two master theorems take a bundle of phases with their proofs, \lean{Phases.Complete} and
\lean{Phases.Security}. Whether these can be met, and by what, is the question.

\paragraph{They can be met by the real phases.} For the bus phase, the table sumcheck and the
public-input phase the review worked out, on paper, a knowledge state function and per-challenge
errors that meet \lean{Phase.Security} against the spine's seams (\cref{ch:phases}); the
public-input phase's is built and proved. For Flock, perfect completeness is attainable by the
honest prover of the specification, and the blueprint's belief that it is not (its acceptance
test 20) is mistaken (\cref{sec:faith-flock}). For the table sumcheck the answer is
\enquote{yes, by a generalized phase with guards leanVM does not have}, which is the finding of
\cref{sec:faith-table}. For the opening, see \cref{sec:faith-opening}.

\paragraph{They can also be met by protocols that are not leanVM's.} Three constructions,
two of them machine-checked:
\begin{enumerate}
\item \textbf{Phases that check nothing, at declared error one.} A phase that draws one
  challenge, ignores it, makes no query and outputs a function of its input statement is
  perfectly complete whenever that function carries one seam into the next, and is
  round-by-round knowledge sound against \emph{any} two seams at the error it declares for its
  challenge, namely $1$. Five of them inhabit \lean{Phases.Security\ I} for every instance
  $I$; both master theorems hold of the bundle; the composed error is $1$ at each of its five
  challenges (probe \code{P4Junk}, standard axioms only). The knowledge-soundness master theorem
  therefore has content only together with a bound on \lean{piopError}, and the spine states
  none: \lean{piopError\ P} is \lean{P.toDef.err}, whatever the phases declare. The blueprint's
  \lean{piopError_le} (Layer 10) is stated over sizes the spine's \lean{piopError} does not take.
  This is the review's first major finding on the spine (\cref{sec:faith-spine}); the proposed
  repair is a closed-form error fixed by the spine from the instance and demanded of each phase.
\item \textbf{Phases that read the whole stack.} In the ideal model the verifier may query
  $\widetilde{q}$ at every point of the cube, and neither ArkLib nor the spine counts queries;
  \lean{M3Holds} is decidable. So a phase with no round whose verifier reads the $2^{\mu}$ cells
  and decides the input seam itself is perfectly complete and knowledge sound at error zero.
  Five of them inhabit \lean{Phases.Security\ I} for every $I$ with \lean{piopError} $= 0$. On
  paper (two dossiers reach it independently); a probe was designed and not run.
\item \textbf{A verifier that always rejects} is knowledge sound at error zero and is excluded
  only because \lean{Security} extends \lean{Complete}: on an instance whose relation is empty
  (below) the exclusion is void.
\end{enumerate}
Conversely, the five pass-through phases of the repository's tests, which inhabit
\lean{Phases.Complete}, have \emph{no} \lean{Phases.Security}: the type is empty (probe
\code{P5PassThrough}, theorem \code{no\_security}), where the test's docstring only asserts it.

\paragraph{What this means.} The master theorems are composition theorems. They pin the
phases to leanVM's by nothing in their statement; what does is outside the spine and, today,
outside the repository: a bound on the error (item 1), the phases' definitions being leanVM's,
which only the compiled verifier's refinement theorem and the fixture of a Rust-produced proof
can tie down (item 2), and review of each phase's definition against the specification. The
report's recommendation is to move the first into the spine and to say the other two plainly
in the blueprint's headline (\cref{ch:options}).

\section{The instance}\label{sec:nv-instance}

The relation is stated over an abstract instance, and an instance can be wrong in ways the
spine does not exclude. A \lean{Layout} that reads every column from cells $0$ and $1$
satisfies the reading law (probe \code{P2Relation}, \code{toyAlias}); on it the three columns
are one, \lean{M3Holds} asks for incompatible things, and the relation is empty for every stack
and statement (exhaustive on the cells, the general case on paper). The master theorems hold of
that instance: completeness is vacuous, and knowledge soundness says that every accepting
prover was lucky. The reading law also determines \lean{read} from \lean{extend} (a cube point's
extension is the cell), so \lean{read} is redundant data and the only freedom of an instance is
\lean{extend}. A wrong \lean{extend} costs completeness and faithfulness, never the soundness of
the Lean chain: the adaptor's theorem from \lean{M3Holds} to \lean{SatisfiedBy}, if proved of
such an instance, still gives the chain to the semantics. Who rules such instances out is the
adaptor's completeness theorems and the transcription of the layout with the compiled verifier's
fixture, none of which is built; until then, no statement on \code{main} says that the relation
of any instance other than the toy is inhabited.

The same holds of the instance's other fields: \lean{aux} is an arbitrary predicate, and on the
only built instance it is \lean{True}, so a Flock phase with no message, no challenge and no
check is knowledge sound at error zero (probe \code{NoCheckFlock}); \lean{publicLines} is a
list, and an instance with two lines instead of three has every theorem of the spine and of the
public-input phase (\cref{ch:drift}); \lean{Layout.comap}, by which the adaptor is to name the
blocks, accepts a renaming that is not one-to-one, so columns can alias.

\section{The error bound}\label{sec:nv-error}

Beyond item 1 above: the errors the blueprint charges to the bus phase, the table sumcheck and
the public-input phase are correct and tight where they are stated, and the review's
round-by-round arguments give the state functions that attain them (\cref{ch:phases}). Two
errors are not tight: the zerocheck escape, which costs $\tau_{\max} / \abs{\E}$ once and
nothing in a conjunctive state function, is charged per constraint and placed on the wrong
challenges; and the GKR round error is $4 / \abs{\E}$, not $5 / \abs{\E}$. Whether the sum
means something at leanVM's parameters is settled in \cref{sec:faith-opening}: the interactive
sum does (about $2^{32.3} / \abs{\E}$ under the stacking window, within the blueprint's
$2^{40} / \abs{\E}$), but the non-interactive
bound the blueprint sketches leaves out three things that decide whether leanVM reaches the
128 bits it claims.

\section{The extractor}\label{sec:nv-extractor}

Every extractor in the repository is a compiled definition: nothing along the chain from
\lean{piopExtractor} to the commit phase's is marked \lean{noncomputable}, and the three of them
run (probes \code{P1Axioms}, \code{Extractors}). \lean{#print\ axioms} says nothing about this:
it reports the three standard axioms for computable definitions whose proofs use classical
lemmas, and the earlier review's use of it as evidence of computability was a misreading.

Two things are less than the blueprint says. First, \enquote{the extractor's column $q$
satisfies \lean{M3Holds}} names no definition of the repository: the composed extractor's
\lean{extractOut} returns the stack only when every phase has no round; with the repository's
own public-input phase in the bundle, the statement of the test that \enquote{the extractor
returns the committed stack} is ill-typed, because \lean{extractOut} then lives in that phase's
last intermediate witness type, \lean{Unit} (probe \code{P3cExtractor}, an expected failure
recorded). The stack is what the composed extractor returns at round $0$ through its
\lean{extractMid} at the commit round, and there it is the transcript's first message whatever
it is handed. No function \enquote{the stack extracted from a full transcript} is defined.
Second, ArkLib's extractor type accepts an extractor that chooses a witness by classical choice,
and such an extractor is \enquote{knowledge sound at error zero} for any relation whose
language the verifier decides (probe \code{Extractors}): ArkLib's existential notion of
round-by-round knowledge soundness is soundness, and the named form the spine uses is knowledge
only with the named extractor read. A phase's \lean{Security} could carry a classical extractor
without anything in the repository noticing; it would not matter, because from the commit phase
on the witnesses are \lean{Unit} and the stack comes from the commit phase's extractor alone,
which is the spine's and reads the message. The design \enquote{the oracle is the witness}
excludes the risk that acceptance test 24 guards against; the test does not.

\section{The public-input phase}\label{sec:nv-pubinput}

The phase is built and its two theorems are proved; the question is what they exclude. The
answer, machine-checked (probe \code{ProbeAnyCheck}, theorem \code{securityG}): nothing about
the check. For every function the prover may use to compute its message and every Boolean check
that the honest message passes, the phase with that check has a \lean{Phase.Security} against
the same seams at the same error $1 / \abs{\E}$. Instances proved: no check and no message; no
check and a message of forty-two zeros; a check with the two cells swapped, with a prover that
sends the swapped values. The reason is one line of the definition: \lean{pooled} takes its
values from the statement, not from the message, and the proof of knowledge soundness never
mentions \lean{check}. What the theorems do exclude, each refuted in Lean with no extractor,
state function or error below one able to repair it (theorem \code{not\_rbr}, which denies
ArkLib's existential notion): a verifier that pools the prover's values without checking them;
one that checks only the first value and pools the prover's values; one that does not pool the
claim of a line whose value is not sent. The earlier review's statement that each of five wrong
verifiers breaks a stated theorem was true of the build it read and is false of the code as
merged: its own compression (pooling the lines' values, its own first compression) removed the
dependence, and the status file still repeats the claim. The consequence for the requirement
that a missing check leave a theorem unprovable is drawn in \cref{ch:drift}; the repair, pooling
the values sent, is proposed in \cref{sec:faith-pub}.

\section{The notion itself}\label{sec:nv-notion}

ArkLib's round-by-round knowledge soundness is not trivially satisfiable: a verifier that
accepts everything has no knowledge state function for a non-trivial relation with no rounds,
and with one challenge it is not knowledge sound at error zero whatever the extractor; every
verifier is knowledge sound at error one (probe \code{NonVacuity}). Its plain reading, that an
accepting transcript either extracts a witness or contains a challenge that fell in the bad set
of its prefix, is machine-checked at the transcript level (probe \code{PlainReading}); the
step from there to a bound on a prover's success is ArkLib's admitted implication from
round-by-round to plain knowledge soundness, at both pins (\cref{ch:libraries}).

\section{The blueprint's acceptance tests}\label{sec:nv-acceptance}

The blueprint names twenty-eight acceptance tests, each a reading that compiles and is wrong
with the witness that rejects it. For each, whether the witness exists on \code{main} and the
review's verdict.
\InputIfFileExists{sections/gen/acceptance-tests}{}{\emph{[table pending]}}

\section{What was not established}\label{sec:nv-open}

The read-everything phase (item 2 of \cref{sec:nv-hypotheses}) is on paper. The emptiness of
the aliased instance's relation is exhaustive on the cells and on paper in general. Of the
mutation probes of the public-input phase, two were never written (the deployed check with the
prover's values pooled, the wrong point) and rest on paper; two (the wrong check with the
honest prover unchanged, the extra check) ran at the new pins and refuted completeness as
predicted, while their unchanged knowledge-soundness proof failed at the state function like
every other mutation's (\cref{ch:probes}).
